% ID: SEL-D0020-OB-E04
% ESTADO: cotejada. Segunda pasada visual completa y borrador compilado revisados.
% Fuente: Selectividad/matematicas_aplicadas/reserva b/Examen-Matemáticas Aplicadas a las Ciencias Sociales.pdf
% SHA-256: 37584a9604930b74008a6bda88ef3d61fa926d5e6a540ab9013569258f772256
% Página: 2; opción B; ejercicio original 4.

La Delegación de Tráfico de una ciudad desea estudiar la influencia del uso del teléfono móvil en los accidentes de tráfico. Elegida una muestra aleatoria simple de 250 accidentes registrados el año pasado, se observó que 90 de ellos se produjeron por distracciones debidas al uso del móvil.
\begin{enumerate}[label=\alph*)]
\item \textbf{(1.5 puntos)} Determine un intervalo de confianza al 97\% para estimar la proporción de accidentes de tráfico debidos al uso del móvil mientras se conduce.
\item \textbf{(1 punto)} Usando la estimación anterior, calcule el tamaño mínimo que debe tener una muestra para estimar la proporción de accidentes con un error máximo del 5\% y un nivel de confianza del 99\%.
\end{enumerate}
